\documentclass[11pt]{article}

\usepackage{ulem}
\usepackage{tikz}
\usepackage{xurl}
\usepackage{float}
\usepackage{amsmath}
\usepackage{booktabs}
\usepackage{graphicx}
\usepackage{listings}
\usepackage[T5]{fontenc}
\usepackage{indentfirst}
\usepackage[utf8]{inputenc}
\usepackage[margin=1in]{geometry}
\usetikzlibrary{arrows.meta, positioning}

\begin{document}

\begin{center}

{\Large \textbf{UNIVERSITY OF SCIENCE AND TECHNOLOGY OF HANOI}}\\
\vspace{0.3cm}
----------------------o0o----------------------\\

\vspace{4cm}

{\large
\centerline{2540027 - Nguyễn Xuân Vinh}
}

\vspace{4cm}

{\Huge \textbf{BC3.04a Advanced Programming for HPC}}\\
\vspace{0.5cm}
Lecturer: Dr. Trần Giang Sơn

\vspace{4cm}

{\Large \textbf{REPORT LABWORK 9}}\\

\vspace{0.5cm}

{\large
Subjects: Histogram Equalization
}

\vspace{4cm}
\today
\end{center}

\newpage

\newpage

\tableofcontents

\newpage

\section{Objective}

The objective of the labwork 9 is to implement histogram calculation and
histogram equalization for a grayscale image.

The lab consists of two main parts:

\begin{itemize}
    \item Labwork 9a: calculate the grayscale histogram using local
    histograms and reduction.
    \item Labwork 9b: perform histogram equalization using the histogram,
    probability distribution, cumulative distribution function (CDF), and
    a pixel intensity mapping.
\end{itemize}

The implementation also tests different 2D block sizes and measures the
execution time of the main operations.

\section{Input Image}

The input image was loaded from:

\begin{verbatim}
/content/img.jpg
\end{verbatim}

The image has a resolution of $1080 \times 1920$ pixels, giving a total of

\[
1080 \times 1920 = 2\,073\,600
\]

pixels.

The input RGB image is first converted to grayscale using OpenCV. The
grayscale image contains intensity values from 0 to 255.

\begin{figure}[H]
    \centering
    \includegraphics[width=0.8\textwidth]{figures/original_grayscale.png}
    \caption{Input grayscale image.}
    \label{fig:original_grayscale}
\end{figure}

\section{Labwork 9a: Histogram}

\section{Histogram Calculation}

A grayscale histogram contains 256 bins. Bin $i$ represents the number of
pixels whose grayscale intensity is equal to $i$.

In a serial implementation, the histogram can be calculated by visiting
every pixel and increasing the corresponding bin:

\begin{verbatim}
for y in range(height):
    for x in range(width):
        histogram[image[y][x]] += 1
\end{verbatim}

For the labwork 9, the histogram is implemented using the local histogram
approach described in the lab sheet.

\section{Local Histograms}

The grayscale image is divided into 2D blocks. Each block calculates its
own local histogram containing 256 bins.

For each pixel in a block, its grayscale value is used as the index of the
corresponding local histogram bin.

Therefore, instead of directly modifying one global histogram, each block
produces an independent local histogram.

\section{Reduction}

After all local histograms have been calculated, they are combined using a
reduction operation.

For every intensity $i$, the global histogram is calculated as

\[
H[i] = \sum_k H_k[i],
\]

where $H_k[i]$ is the local histogram of block $k$.

This produces the final 256-bin histogram of the complete image.

For the main experiment, a block size of $16 \times 16$ was used.

The calculated histogram contained exactly

\[
2\,073\,600
\]

pixels, which is equal to the total number of pixels in the input image.
The histogram was also compared with a serial reference implementation and
the result was identical.

\section{Labwork 9b: Histogram Equalization}

Histogram equalization increases the global contrast of an image by
remapping the original grayscale intensities according to the cumulative
distribution of the histogram.

The implementation follows the steps given in the lab sheet.

\section{Probability Distribution}

First, the probability of each intensity value is calculated from the
histogram:

\[
p_i = \frac{H[i]}{n},
\]

where $H[i]$ is the histogram count for intensity $i$ and $n$ is the total
number of pixels.

For this image,

\[
n = 2\,073\,600.
\]

The sum of all probabilities was

\[
\sum_{i=0}^{255} p_i = 1.000000000000.
\]

\section{Cumulative Distribution Function}

The cumulative distribution function is calculated sequentially:

\[
c_i = \sum_{j=0}^{i} p_j.
\]

Although a parallel Scan could be used for this operation, the lab sheet
specifically states that Scan has not yet been introduced. Therefore, the
CDF was calculated sequentially.

The final CDF value was

\[
c_{255} = 1.000000000000.
\]

The calculated CDF was also verified to be monotonically increasing and to
remain within the range $[0,1]$.

\section{Intensity Mapping}

The CDF is linearly scaled from $[0,1]$ to the grayscale range $[0,255]$:

\[
h_i = c_i \times 255.
\]

The resulting mapping has a range from 0 to 255.

Each original grayscale pixel with intensity $i$ is then replaced by
$h_i$.

This produces the histogram-equalized image.

\begin{figure}[H]
    \centering
    \includegraphics[width=0.8\textwidth]{figures/equalized.png}
    \caption{Histogram-equalized grayscale image.}
    \label{fig:equalized}
\end{figure}

\section{Results}

\section{Correctness}

Several checks were performed after the implementation.

The histogram contained exactly the same number of pixels as the input
image:

\[
\sum_{i=0}^{255} H[i] = 2\,073\,600.
\]

The probability distribution summed to 1, and the final CDF value was also
1.

The following correctness checks were obtained:

\begin{itemize}
    \item Histogram count correct: True.
    \item Probability sum approximately 1: True.
    \item CDF monotonic: True.
    \item CDF is in $[0,1]$: True.
    \item Equalization mapping is monotonic: True.
    \item Equalized image size unchanged: True.
    \item Equalized histogram count correct: True.
\end{itemize}

\section{Different Block Sizes}

The implementation was tested with five different 2D block sizes.

\begin{table}[H]
\centering
\caption{Performance for different block sizes.}
\label{tab:block_results}
\begin{tabular}{lrrrr}
\toprule
Block Size & Histogram (s) & CDF (s) & Equalization (s) & Total (s) \\
\midrule
$4\times4$   & 1.763331 & 0.000074 & 0.305602 & 2.069006 \\
$8\times8$   & 1.631795 & 0.000130 & 0.142211 & 1.774136 \\
$16\times16$ & 1.378689 & 0.000065 & 0.024245 & 1.403000 \\
$32\times32$ & 0.811796 & 0.000115 & 0.010864 & 0.822775 \\
$64\times64$ & 0.857225 & 0.000106 & 0.007107 & 0.864438 \\
\bottomrule
\end{tabular}
\end{table}

All block sizes produced a correct histogram. The equalized images produced
with the different block sizes were also identical to the $16\times16$
reference result.

The maximum difference between the results was

\[
0
\]

and the mean difference was

\[
0.0.
\]

Therefore, changing the block size affected the execution time but did not
change the resulting histogram or equalized image.

\section{Performance}

The execution time was measured for the histogram calculation, the
sequential CDF calculation, and the equalization mapping.

The histogram calculation was the main computational part of the
implementation. Its execution time decreased when the block size was
increased from $4\times4$ to $32\times32$. The $64\times64$ configuration
had a slightly higher histogram time than $32\times32$ in this experiment.

The equalization mapping also became faster with larger blocks, decreasing
from $0.305602$ seconds for $4\times4$ blocks to $0.007107$ seconds for
$64\times64$ blocks.

The total measured time ranged from $2.069006$ seconds for $4\times4$
blocks to $0.864438$ seconds for $64\times64$ blocks.

No conventional speedup value was calculated because the lab implementation
does not contain a separate optimized implementation to use as a baseline.
The measurements therefore show the performance differences between block
sizes rather than a speedup between two different algorithms.

\section{Conclusion}

In the labwork 9, a 256-bin histogram was implemented for a grayscale image
using local histograms followed by a reduction operation. The histogram
calculation was verified against a serial reference implementation.

Histogram equalization was then implemented using the histogram,
probability distribution, sequential CDF calculation, and intensity
mapping from $[0,255]$ to $[0,255]$.

The implementation was tested with block sizes of $4\times4$, $8\times8$,
$16\times16$, $32\times32$, and $64\times64$. All configurations produced
the same correct histogram and identical equalized output.

The experiments also showed that the choice of block size affects
execution time, while the resulting image remains unchanged.

\end{document}